\section{\texorpdfstring{$\epsilon$}{ε}-正则性与密度估计的证明}
\label{sec:epsilon-regularity-density-estimates}

内部主能量密度估计（定理 8.1.3）本质地使用逐点微分不等式
$\Delta e(u)\leq C_1e(u)+C_2e(u)^2$ 与平均值定理；这种局部论证不能直接用于边界估计。因此，我们采用 (Oh92) 的更整体的方法，在同一个框架下得到边界估计和内部估计。这个策略又在很大程度上把萨克斯--乌伦贝克 (SU81) 研究调和映射时的整体方法移植到当前的伪全纯背景。特别地，我们把他们的方案用于边值问题。小能量 $J$-全纯映射先验估计的关键，是把辛几何与 (SU81) 中著名的萨克斯--乌伦贝克重标度论证结合起来。这简化了边界估计的研究；否则便须像叶如钢在 (Ye94) 中所作的那样，对自由边值问题作一番非平凡的研究。

始终假设 $M$ 闭，或在非紧时具有有界几何。这一节要证明的主要定理如下。以
$\mathring D$ 表示开单位圆盘或以 $(-1,1)\subset\mathbb C$ 为边界的开半圆盘（见图 8.1）。

\begin{theorem}\label{thm:ch8-epsilon-regularity-w1p}
设 $R\subset M$ 为全实子流形，$u:D\to M$ 满足 $u(\partial D)\subset R$ 且
$\bar\partial_Ju=0$。则存在 $\epsilon>0$，使得若
$\int_{D(1)}|du|^2<\epsilon$，则对任意 $p\geq2$ 及任意
$\overline{D'}\subset\mathring D$，有 $\|du\|_{p,D'}\leq K_4$；其中
$K_4$ 依赖于 $\epsilon,p,J,D'$。
\end{theorem}

由此立即可得下述导数界；在几何测度论的文献中，它常被称为 $\epsilon$-正则性定理。
\begin{corollary}[$\epsilon$-正则性定理]\label{cor:ch8-epsilon-regularity}
\stfhIndex{I041}{00-symbol-041}{$\epsilon$-正则性定理（$\epsilon$-regularity theorem）}
设 $\overline{D'}\subset\mathring D$，并令 $u$ 满足定理 8.4.1 的条件，则
$\|du\|_{\infty,D'}\leq K_5$。
\end{corollary}
\begin{proof}
取 $\overline{D'}\subset D''$ 且 $\overline{D''}\subset\mathring D$。对 $(\overline{D'},D'')$ 应用定理 8.4.1，得
$\|du\|_{1,p,\overline{D'}}\leq K_6$；索伯列夫不等式给出
$\|du\|_{\infty,D'}\leq C\|du\|_{1,p,D''}$。
\end{proof}

\begin{remark}\label{rem:ch8-rescaling-lp}
若 $u:\Sigma\to M$ 光滑，$B_\eta(x)\subset\Sigma$，定义
$u_\eta:D(1)\to M$ 为 $u_\eta(z)=u(x+\eta z)$，则
\[
 \int_{D(1)}|du_\eta|^p=\eta^{p-2}\int_{B_\eta(x)}|du|^p,\qquad
 \|du_\eta\|_{p,D(1)}=\eta^{1-2/p}\|du\|_{p,B_\eta(x)}.
\]
当 $p=2$ 时重标度不改变导数的 $L^p$ 范数；当 $p>2$ 时，右端随
$\eta\to0$ 而趋于零。
\end{remark}

\begin{proof}[定理 8.4.1 的证明]
反设存在 $\overline{D'}\subset\mathring D$ 及序列 $u_\alpha$，满足边界与全纯条件且
\begin{equation}
 \int_{D(1)}|du_\alpha|^2\to0,\qquad
 \|du_\alpha\|_{p,D'}\to\infty
 \label{eq:ch8-epsilon-regularity-contradiction}
\end{equation}
当 $\alpha\to\infty$。定义
\[
 \eta_\alpha(D')=\inf\{\eta>0\mid\exists x\in\overline{D'}:
 \eta^{1-2/p}\|du_\alpha\|_{p,B_\eta(x)}\geq1\}.
\]
对每个固定 $\alpha$，此下确界严格为正；否则可取
$\eta_j\to0$ 与 $x_j\in\overline{D'}$，使局部范数至少为
$\eta_j^{-1+2/p}\to\infty$，与 $\|du_\alpha\|_{p,D}<\infty$ 矛盾。
现在让 $\alpha$ 变化。由 $\eta_\alpha(D')>0$ 的定义，对所有
$x\in\overline{D'}$ 以及 $\eta<\eta_\alpha(D')$，有
\[
 \eta^{2-p}\|du_\alpha\|_{p,B_\eta(x)}^2\leq1.
\]
若 $\eta_\alpha(D')$ 在某子列上远离零，用有限个半径
$\eta_0$ 的闭球覆盖 $D'$。可取
\[
0<\eta_0<\min\{\eta_\alpha(D'),\operatorname{dist}(\overline D-D')\},
\]
并用与 $\alpha$ 无关的 $N_{\eta_0}$ 个闭球
$B_{\eta_0}(x_{\alpha,k})\subset D$（$k=1,\ldots,N_{\eta_0}$）覆盖 $D'$，
使得 $x_{\alpha,k}\in\overline{D'}$；在每个球上都有
\[
\eta_0^{1-2/p}\|du_\alpha\|_{p,B_{\eta_0}(x_{\alpha,k})}\leq1.
\]
于是
\[
 \|du_\alpha\|_{p,B_{\eta_0}(x_{\alpha,k})}^p\leq\eta_0^{2-p}.
\]
把此不等式对这些闭球相加，再取“$p$ 次平方根”，便得到
\[
 \|du_\alpha\|_{p,D'}\leq C N_{\eta_0}\eta_0^{2/p-1},
\]
这与
\eqref{eq:ch8-epsilon-regularity-contradiction} 矛盾。因此
$\eta_\alpha(D')\to0$ 当 $\alpha\to0$。利用 $\eta_\alpha(D')$ 的定义、
$\eta^{2-p}\|du_\alpha\|_{p,B_\eta(x)}^p$ 关于 $\eta$ 的连续性以及上述评注，
可对每个给定 $\alpha$ 把 $\eta_\alpha$ 选得任意接近但小于
$\eta_\alpha(D')$，并取
$x_\alpha\in\overline{D'}$，使
\begin{equation}
 \int_{B_{\eta_\alpha}(x_\alpha)}|du_\alpha|^p
 \geq\frac12\eta_\alpha^{2-p},
 \label{eq:ch8-rescaling-lower-bound}
\end{equation}
且对所有 $x\in\overline{D'}$，
\begin{equation}
 \eta_\alpha^{p-2}\|du_\alpha\|_{p,B_{\eta_\alpha}(x)}^2\leq1.
 \label{eq:ch8-rescaling-upper-bound}
\end{equation}
令
$r_\alpha=\eta_\alpha^{-1}\min\{\operatorname{dist}(x_\alpha,\partial D),
\operatorname{im}(x_\alpha)\}$。取子列后有两种情形：
$r_\alpha\to\infty$，或 $r_\alpha\to r<\infty$。

第一种情形定义 $v_\alpha(z)=u_\alpha(\eta_\alpha z+x_\alpha)$。其定义域是
\[
 \{z\in\mathbb C\mid\eta_\alpha z+x_\alpha\in D\},
\]
这逐项等价于
\[
 \begin{cases}|\eta_\alpha z+x_\alpha|\leq1,\\
 \operatorname{im}(\eta_\alpha z+x_\alpha)\geq0
 \end{cases}
 \Longleftrightarrow
 \begin{cases}|z+x_\alpha/\eta_\alpha|\leq1/\eta_\alpha\to\infty,\\
 \operatorname{im}z\geq-\operatorname{im}(x_\alpha/\eta_\alpha).
 \end{cases}
\]
由于 $1-|x_\alpha|\geq\eta_0>0$ 且
$\operatorname{im}(x_\alpha/\eta_\alpha)\to\infty$，有
$1/\eta_\alpha-|x_\alpha|/\eta_\alpha=(1-|x_\alpha|)/\eta_\alpha\to\infty$，
故其定义域最终包含任意
$B_R(0)$，并可假设 $B_R(0)$ 包含在定义域内。于是映射
$v_\alpha:B_R(0)\to M$ 满足
\begin{enumerate}[label=(\arabic*)]
\item $\|dv_\alpha\|_2\to0$；
\item $\|dv_\alpha\|_{p,B_1(0)}^p\geq\frac12$\footnote{译者注：原书此式无上标 $p$，但按 (8.4.29) 应为 $p$ 次幂；译文补出。}；
\item 对 $x\in B_{R-1}(0)$，$\|dv_\alpha\|_{p,B_1(x)}\leq1$；
\item $\bar\partial_Jv_\alpha=0$。
\end{enumerate}
这里 (3) 来自 \eqref{eq:ch8-rescaling-upper-bound}。对每个固定的 $R$，考虑
$v_\alpha|_{B_R}$ 的极限。由 (3) 及主要局部估计，
$\|dv_\alpha\|_{1,p;B_{9/10}(x)}\leq C(R)$。索伯列夫嵌入给出一个子列，使它对某个 $\epsilon>0$ 在每个
$B_{8/10}(x)$ 上的赫尔德空间 $C^\epsilon$ 中收敛；再应用定理 8.3.5 及 bootstrap，可取子列在每个 $B(R)$ 上
$C^\infty$ 收敛到 $w_R$。极限 $w_R:B_R(0)\to M$ 满足：
\begin{enumerate}[label=(\arabic*)]
\item $\bar\partial_Jw_R=0$；
\item $\|dw_R\|_2^2\leq\limsup_\alpha\|dv_\alpha\|_{2;B_R}^2\to0$，因而
$\|dw_R\|_2^2=0$；
\item 因 $v_\alpha\to w_R$ 于 $C^1$，
$\|dw_R\|_{p,B_1(0)}^p=\lim_{\alpha\to\infty}
\|dv_\alpha\|_{p,B_1(0)}^p\geq1/2$。
\end{enumerate}
(2) 与 (3) 显然矛盾，故第一种情形不可能。

下面考虑第二种情形，即 $r_\alpha\to r<\infty$。由于假设
$D'\subset\overline{D'}\subset D$，特别是在半圆盘情形，最终这个最小值由
$\eta_\alpha^{-1}\operatorname{im}(x_\alpha)$ 取得。在这种情形，
$x_\alpha\to\partial D$ 与 $\eta_\alpha\to0$ 一样快。
\begin{figure}[htbp]
\centering
\begin{tikzpicture}[scale=.8]
 \draw (-2.7,.25)--(2.7,.25);
 \draw (-2.3,.25) arc[start angle=180,end angle=0,radius=2.3];
 \draw (-1.85,-.35)--(1.85,-.35);
 \draw[dashed,->] (.5,1.35)--(.5,.38);
 \node[right] at (.65,1.05) {$x_\alpha$};
 \node[below] at (0,-.35) {$\partial D$};
\end{tikzpicture}
\caption{趋向边界的点。}
\label{fig:ch8-points-tending-boundary}
\end{figure}
此时同样有
\[
 \begin{cases}|\eta_\alpha z+x_\alpha|<1,\\
 \operatorname{im}(\eta_\alpha z+x_\alpha)\geq0
 \end{cases}
 \Longleftrightarrow
 \begin{cases}|z+x_\alpha/\eta_\alpha|<1/\eta_\alpha\to\infty,\\
 \operatorname{im}z\geq-\operatorname{im}(x_\alpha/\eta_\alpha)\to-r.
 \end{cases}
\]
重标度定义域最终包含
$\{z\in\mathbb C\mid |z|\leq R,\operatorname{im}z\geq-r\}$。采用同样的取极限过程，仍会得到矛盾，故第二种情形也不可能。这就证明了定理 8.4.1 中所需的 $\epsilon>0$，从而完成证明。
\end{proof}

下面证明如下的一致估计。
\begin{proposition}\label{prop:ch8-uniform-derivative-estimate}
设 $u,\epsilon$ 如 $\epsilon$-正则性定理（即 $\int|\nabla u|^2<\epsilon$ 且
$\bar\partial_Ju=0$）。若 $0<r<r'<1$，则
$\max_{|x|<r}|du(x)|\leq CE(u|_{D(r')})$，其中 $C=C(r,r')$ 与 $u$ 无关，只依赖于 $r,\epsilon$。
\end{proposition}
\begin{proof}
在 $D(r')$ 上对 $K=D(r)$ 应用主要椭圆估计；这里 $D$ 并非单位圆盘，但重标度后得到同类估计，其常数只依赖于 $r$。又因
$E(u|_{D(r')})\leq E(u)\leq\epsilon$，$\epsilon$-正则性给出与 $u$ 无关的
$|du|\leq K_5$。命题 8.3.4 因而化为
\[
 \int_{D(r)}|\nabla(du)|^2\leq C\int_{D(r')}|du|^2(1+|du|+|du|^2)
 \leq C\int_{D(r')}(1+K_5+K_5^2)|du|^2.
\]
令 $C'=C(1+K_5+K_5^2)$，对 $r<r''<r'$ 得
$\|du\|_{1,2;D(r'')}\leq C\|du\|_{2,D(r')}$。再由索伯列夫不等式，
\[
 \max_{x\in D(r)}|du(x)|\leq K'\|du\|_{1,2;D(r'')}
 \leq K'K(r)\|du\|_{2,D(r')}.
\]
令 $C=K'K(r)$，并注意 $\|du\|_{2,D(r')}=E(u|_{D(r')})$，便完成了证明。
\end{proof}

\begin{remark}\label{rem:ch8-rectifies-oh92}
上述证明纠正了 (Oh92) 中相应命题证明的一处错误。
\end{remark}

下面这个关键估计是定理 8.1.3 的边界版本。
\begin{proposition}[边界能量密度估计]\label{prop:ch8-boundary-energy-density}
\stfhSubIndex{I117}{neng^liang^mi^du^gu^ji}{能量密度估计（energy density estimate）}{bian^jie}{边界（boundary）}
设 $U,\epsilon$ 如上，但假设 $U$ 定义在 $D(2)$ 上。则对 $x\in D(1/2)$，
\[
 |x||du(x)|\leq\sqrt C\left(\int_{D(2|x|)}|du(x)|^2\right)^{1/2},
\]
等价地，若 $r=|x|$，则 $|du(x)|^2\leq CE(2r)/r^2$。
\end{proposition}

\begin{remark}\label{rem:ch8-boundary-mean-value}
令 $v(z)=u(|x|z)$，则
\begin{equation}
 |dv(0)|\leq C\int_{D(2)}|dv(z)|^2.
 \label{eq:ch8-boundary-mean-value}
\end{equation}
这是平均值不等式 \eqref{eq:ch8-morrey-mean-value} 的边界版本；该命题的证明也适用于定理 8.1.3 的内部估计。
\end{remark}
\begin{proof}
取 $x_0\in D(1/2)$，在
$N(3/2,x_0)$ 上考虑 $v(z)=u(x_0+|x_0|z)$，其中
\[
 N(r,x_0)=\{z\in D(r)\mid\operatorname{Im}z\geq-\operatorname{Im}x_0/x_0\}.
\]
若 $x_0\in D(1/2)$，则有共形等价
\[
 (N(3/2,x_0),\partial N(3/2,x_0))\sim(D(1),\partial D(1)),
\]
且该共形映射 $\phi$ 的导数对 $x_0\in D(1/2)$ 一致有界。考虑
$v\circ\phi^{-1}:D(1)\to M$；能量共形不变，故命题 8.4.4 给出
$|dv(0)|\leq C\|dv\|_{2,D(1)}$。结合 $dv(0)=|x_0|du(x_0)$ 与
$\|dv\|_{2,D(1)}=\|du\|_{2,N(3/2,x_0)}
\leq\|du\|_{2,D(2|x_0|)}$ 即证。
\end{proof}

事实上，我们可以量化 $\epsilon$-正则性定理中出现的 $\epsilon$ 能有多大。为此引入几何量
\begin{equation}
 A(\omega,L;J)=\min\{A_S(\omega,J),A_D(\omega,J;L)\}>0,
 \label{eq:ch8-minimal-bubbling-energy}
\end{equation}
其中
\[
 A_S(\omega,J):=\inf\left\{\int v^*\omega>0\mid
 v:S^2\to M,\ \bar\partial_Jv=0\right\},
\]
以及
\[
 A_D(\omega,L;J):=\inf\left\{\int w^*\omega>0\mid
 w:(D^2,\partial D^2)\to(M,L),\ \bar\partial_Jw=0\right\}.
\]

\begin{proposition}\label{prop:ch8-positive-bubbling-threshold}
设 $L$ 是 $M$ 的拉格朗日子流形，则：(1) $A(\omega,L;J)>0$；
(2) 对任意 $\delta>0$，定理 8.4.1 对所有满足
$E_J(w)<A(\omega,L;J)-\delta$ 的 $w$ 一致成立，估计常数只依赖于
$\delta,L,J$。
\end{proposition}
\begin{proof}
先作如下准备。达布--温斯坦定理给出微分同胚 $\psi:V\to U$，其中
$V$ 是 $T^*L$ 零截面的邻域，$U$ 是 $M$ 中 $L$ 的邻域，并满足
$\psi^*\omega=\omega_0$、$\omega_0=-d\theta$。这里 $\theta$ 是 $T^*L$ 的典范
$1$-形式，按定义在零截面上恒为零。若 (1) 不成立，则存在序列 $u_\alpha$，满足
$u_\alpha|_{\partial\Sigma}\subset L$、$E(u_\alpha)>0$、
$E(u_\alpha)\to0$、$\bar\partial_Ju_\alpha=0$。由命题 8.4.4，导数一致趋零，像最终落在 $L$ 的达布--温斯坦邻域。于是
$0<\frac12E(u_\alpha)=\int u_\alpha^*\omega=-\int_{\partial\Sigma}u_\alpha^*\theta=0$，矛盾。(2) 只要对任意给定的 $\delta>0$ 假设
$E(u_\alpha)<A(\omega,L;J)-\delta$，同样仔细检查上述证明也会导出矛盾，细节留作习题。
\end{proof}

\begin{exercise}\label{ex:ch8-uniform-bubbling-threshold}
补全上面陈述 (2) 的证明。
\end{exercise}

这一命题引出一个有趣辛不变量的如下定义；切卡诺夫 (Che98) 最先使用了它。
关于这种 $\epsilon$-正则性型不变量的用途，亦见 (Oh97c)、(Oh05d)。
\begin{definition}\label{def:ch8-lagrangian-epsilon-invariant}
设 $(M,\omega)$ 驯服，$\mathcal J_\omega$ 为驯服殆复结构集合。对无边界紧拉格朗日子流形 $L$，定义
\[
 A(L;M,\omega):=\sup_{J\in\mathcal J_\omega}A(\omega,J;L).
\]
\end{definition}
若存在 $J\in\mathcal J_\omega$，使得既无 $J$-全纯球面又无边界在 $L$ 上的
$J$-全纯圆盘，则 $A(L;M,\omega)=\infty$；弱正合的 $L$ 即为一例。

\begin{corollary}\label{cor:ch8-weakly-exact-c0-bound}
设 $M$ 闭且 $L$ 弱正合。固定带边黎曼曲面 $(\Sigma,j)$，考虑有限能量的
$(j,J)$-全纯映射 $u:(\Sigma,\partial\Sigma)\to(M,L)$。若
$E_{(j,J)}(u)<\infty$，则存在只依赖于 $E$ 的 $C>0$，使所有
$E_{(j,J)}(u)\leq E$ 的映射满足 $|du|_{C^0}\leq C$。
\end{corollary}

\begin{remark}\label{rem:ch8-rational-lagrangian}
若 $L$ 有理，则命题 8.4.8 显然。若 $L$ 无理，$A(\omega,L;J)$ 的正性没有拓扑来源；此时正性来自几何来源，更确切地说来自 $\epsilon$-正则性结果。在一些文献中，该推论被说成格罗莫夫紧性定理的推论；这是一种误导性说法，因为对无理的 $L\subset(M,\omega)$，格罗莫夫紧性定理本身的证明恰恰依赖此推论。
\end{remark}

现在简要研究一般殆复流形 $(N,J)$ 的紧全实边界条件 $R$ 的情形。这里紧随叶如钢在 (Ye94) 中的论述，只把他的二阶方法稍作修改，换成一阶方法。固定引理 8.3.1 给出的殆埃尔米特度量 $g$，并定义
\[
 A(J,R;g):=\{w:D^2\to N\mid\bar\partial_Jw=0, w(\partial D^2)\subset R, w\text{ 非恒定}\},
\]
\[
 A(J;g):=\{v:S^2\to N\mid\bar\partial_Jv=0, v\text{ 非恒定}\}.
\]
下面的结果是叶如钢 (Ye94) 对带自由边界的调和映射型方程所得一般结论的一个特例。这里采用纯一阶方法，它只涉及带诺伊曼边界条件的伪全纯曲线方程。

\begin{theorem}\label{thm:ch8-positive-energy-totally-real}
若殆埃尔米特流形 $(N,J,g)$ 具有有界几何，则 $A(J;g)>0$；若
$R\subset(N,J)$ 是紧全实子流形，则 $A(J,R;g)>0$。
\end{theorem}
\begin{proof}
只论 $A(J,R;g)$，把 $A(J;g)$ 留给读者。反设 $A(J,R;g)=0$，则有子列
$w_\alpha$ 满足 $w_\alpha|_{\partial\Sigma}\subset R$、$E(w_\alpha)>0$、
$E(w_\alpha)\to0$、$\bar\partial_Jw_\alpha=0$。最终
$E(w_\alpha)<\epsilon$。用有限个单位圆盘和半圆盘覆盖 $\Sigma$，并使半径
$1/2$ 的圆盘已经覆盖 $\Sigma$，则所有 $W^{1,p}$ 范数有一致界。因此
$w_\alpha$ 收敛到常映射 $p\in R$，且其像落在 $p$ 的小邻域。取局部坐标
$\phi:(U,\partial U\cap R)\to(\mathbb C^n,\mathbb R^n)$，满足
$\phi(p)=0$ 且 $\phi_*J=i$ 于 $\mathbb R^n$；于是 $\phi\circ u_\alpha$ 为
$\phi_*J$-全纯。

\begin{exercise}\label{ex:ch8-chart-composition-holomorphic}
证明最后这个陈述。
\end{exercise}

\begin{lemma}\label{lem:ch8-normal-derivative-perpendicular}
若 $\nu$ 是 $\partial\Sigma$ 的外单位法向量，则
$(\partial\phi\circ u_\alpha)/\partial\nu\perp\mathbb R^n\subset\mathbb C^n$。
\end{lemma}
\begin{proof}
在给定共形类 $j$ 中选埃尔米特度量 $h$，并在边界附近取柱坐标
$(r,\theta)$，使 $h=\lambda(dr^2+d\theta^2)$ 且
$j\partial/\partial r=\partial/\partial\theta$。方程等价于
$\partial u/\partial r+J\partial u/\partial\theta=0$。因边界条件使
$\partial u_\alpha/\partial\theta\in TR$，继而有
$J\partial u/\partial r\in TR$。由 $\phi_*J=i$ 于 $\mathbb R^n$ 及局部坐标的选择，故
$\partial(\phi\circ u_\alpha)/\partial r\perp\mathbb R^n$。这就完成了证明。
\end{proof}

定理 8.4.13 的证明由下述习题得到。
\begin{exercise}\label{ex:ch8-subharmonic-coordinate-square}
令 $f_\alpha(z)=|\phi\circ u_\alpha|^2$。证明
$\partial f_\alpha/\partial\nu|_{\partial D^2}=0$；并在 $\phi_*J$ 于像上
$C^2$-接近标准复结构 $i$ 时证明 $\Delta f_\alpha\geq0$。
\end{exercise}
有了这个习题，强最大值原理蕴含 $f_\alpha$ 必为常数。但这与
$E(u_\alpha)>0$（从而 $E(\phi\circ u_\alpha)>0$）矛盾。这就完成了定理 8.4.13 的证明。
\end{proof}

\begin{remark}\label{rem:ch8-no-topological-bound-totally-real}
对与辛形式相容的 $J$ 及拉格朗日边界，映射的整体能量界由引理 7.2.3 的映射拓扑给出；一般全实情形没有这样的拓扑界。
\end{remark}
